% Source: https://artofproblemsolving.com/community/c3h3248729_2024_emcc_guts_round_p1_p12_exeter_math_club_competition
% Stable identity: aops-topic-3248729-q07
An improper fraction of value above one must have a one-digit denominator and at most a two-digit numerator. A denominator of one reduces to an integer, whose decimal representation fits. With denominator two, every allowed numerator gives an integer or half-integer of at most two digits before the decimal point, so the exact decimal uses at most four characters. With denominator at least three, the largest reduced fraction is $\dfrac{98}{3}$: numerator $99$ with denominator three would reduce, and denominators at least four give at most $\dfrac{99}{4}<\dfrac{98}{3}$. Fractions with two-digit denominators and one-digit numerators are below one. Since $\dfrac{98}{3}$ has a repeating nonterminating decimal, the answer is $\boxed{\dfrac{98}{3}}$.
